%

\section{Twierdzenie Desargues'a} % lang-pl
\section{Teorema di Desargues} % lang-it

\begin{proposition}
[twierdzenie Desargues'a] % lang-pl
[teorema di Desargues] % lang-it
\index{twierdzenie!Desargues'a}% % lang-pl
\index{teorema!di Desargues}% % lang-it
    Niech $ABC$ i $A'B'C'$ będą dwoma trójkątami o parami różnych wierzchołkach, przy czym żaden wierzchołek jednego trójkąta nie leży na prostej zawierającej bok drugiego. % lang-pl
    Wtedy proste $AA'$, $BB'$, $CC'$ są współpękowe (mówimy, że trójkąty są \emph{perspektywiczne względem punktu}) wtedy i tylko wtedy, gdy punkty przecięcia odpowiednich boków $AB \cdot A'B'$, $BC \cdot B'C'$, $CA \cdot C'A'$ są współliniowe (mówimy, że trójkąty są \emph{perspektywiczne względem prostej}) -- patrz Neugebauer, strona 109. % lang-pl
    Sia $ABC$ e $A'B'C'$ due triangoli con vertici a due a due distinti, tali che nessun vertice dell'uno giaccia sulla retta contenente un lato dell'altro. % lang-it
    Allora le rette $AA'$, $BB'$, $CC'$ sono concorrenti (diciamo che i triangoli sono \emph{prospettivi rispetto a un punto}) se e solo se i punti di intersezione dei lati corrispondenti $AB \cdot A'B'$, $BC \cdot B'C'$, $CA \cdot C'A'$ sono allineati (diciamo che i triangoli sono \emph{prospettivi rispetto a una retta}) -- si veda Neugebauer, pagina 109. % lang-it
    % TODO: https://en.wikipedia.org/wiki/Desargues%27s_theorem

    Eves s. 245: Desargues two triangle theorem for the plane: copolar triangles are coaxial % lang-pl
    Eves p. 245: teorema dei due triangoli di Desargues nel piano: triangoli copolari sono coassiali % lang-it
\end{proposition}
% 1. Zna pojęcie inwolucji rzutowych.  Zna i potrafi stosować twierdzenia inwolucyjne Desarguesa.

Piszą o nim Audin \cite[s. 26, 151]{audin_2003}, Hartshorne w formie ćwiczenia \cite[s. 183]{hartshorne_2000}. % lang-pl
Eves \cite[s. 362]{eves1_1972} poda przykład Foresta Raya Moultona z 1902 roku płaszczyzny rzutowej, na której nie jest spełnione (!). % lang-pl
Ne parlano Audin \cite[p. 26, 151]{audin_2003} e Hartshorne sotto forma di esercizio \cite[p. 183]{hartshorne_2000}. % lang-it
Eves \cite[p. 362]{eves1_1972} fornirà un esempio di piano proiettivo del 1902, dovuto a Forest Ray Moulton, nel quale esso non è soddisfatto (!). % lang-it
\index[persons]{Moulton, Forest Ray}% https://en.wikipedia.org/wiki/Forest_Ray_Moulton
% Eves 365: Desaurgian/Pappian/classic
Marek Kordos opisze w $\Delta_{15}^{3}$, że z twierdzenia Desargues'a nie wynika twierdzenie Pascala (w wersji dla sześciokąta wpisanego w dwie proste; w płaszczyźnie rzutowej nad kwaternionami zachodzi pierwsze, a drugie nie), natomiast Gerhard Hessenberg udowodni, że z twierdzenia Pascala wynika twierdzenie Desargues'a, i przytoczy ,,grecki'' dowód tego faktu. % lang-pl
Marek Kordos descriverà in $\Delta_{15}^{3}$ che dal teorema di Desargues non segue quello di Pascal (nella versione per l'esagono inscritto in due rette; nel piano proiettivo sui quaternioni vale il primo, ma non il secondo), mentre Gerhard Hessenberg dimostrerà che dal teorema di Pascal segue quello di Desargues, e riporterà una dimostrazione ``alla greca'' di questo fatto. % lang-it
Joanna Jaszuńska pokaże w $\Delta_{15}^{4}$, jak z twierdzenia Desargues'a wyprowadzić twierdzenie o nożycach (o współliniowości punktów przecięcia przekątnych czworokątów wyznaczonych przez pęk prostych i dwie proste go przecinające) i jak wykorzystać oba twierdzenia do konstrukcji wykonywanych samą linijką. % lang-pl
Joanna Jaszuńska mostrerà in $\Delta_{15}^{4}$ come dedurre dal teorema di Desargues il teorema delle forbici (sull'allineamento dei punti di intersezione delle diagonali dei quadrilateri determinati da un fascio di rette e da due rette che lo tagliano) e come usare entrambi i teoremi per costruzioni eseguite con la sola riga. % lang-it

https://www.deltami.edu.pl/2015/04/desargues-i-nozyce/

https://www.deltami.edu.pl/2015/03/dowod-w-stylu-greckim/

%